Let a redundancy-\(r\) generalized Reed--Solomon code have evaluation set
\(S=\mathbf P^1(\mathbb F_q)\setminus A\), where \(|A|=s\), with nonzero
column multipliers.  For \(r\geq6\) and
\[
q\geq6(r+s)-16+\left\lfloor2\sqrt{6(r+s)-18}\right\rfloor,
\]
we classify every coset of weight at least \(r-1\) in odd characteristic and
in characteristic two when \(r\geq8\).  For binary \(r\in\{6,7\}\),
classification is complete for full support; for deleted support we prove
a carrier containment.  If \(s=0\), the covering
radius is \(r-1\), and its maximum-weight directions are the tangent and
conjugate-secant points of the normal rational curve.  If \(s>0\), the radius
is \(r\): omitted curve points have weight \(r\), while tangents, conjugate
secants, and rational secants incident with \(A\) give exactly weight \(r-1\).
We count the shells and their MDS/NMDS extensions and aggregate enumerators.

The proof converts low-weight syndrome representations into split forms in a
Hankel kernel.  A simultaneous contraction reduces arbitrary redundancy to a
cubic terminal problem.  An exact genus-one count constructs witnesses outside
the rank-two locus of the second catalecticant; above the threshold only the
two characteristic-two carriers can obstruct it.  Artifacts record small-field
and trust boundaries.
